\documentclass[10pt,a4paper]{amsart}
\usepackage{amsmath,amssymb,amsthm,graphicx}
% For commutative diagrams we use Xy-pic
\usepackage[all]{xy}

% We use 2cell for 2-commutative diagrams.
\xyoption{2cell}
% Theorem environments.
%
\theoremstyle{plain}
\newtheorem{theorem}[subsection]{Theorem}
\newtheorem{proposition}[subsection]{Proposition}
\newtheorem{lemma}[subsection]{Lemma}

\theoremstyle{definition}
\newtheorem{definition}[subsection]{Definition}
\newtheorem{example}[subsection]{Example}
\newtheorem{exercise}[subsection]{Exercise}
\newtheorem{situation}[subsection]{Situation}

\theoremstyle{remark}
\newtheorem{remark}[subsection]{Remark}
\newtheorem{remarks}[subsection]{Remarks}

\numberwithin{equation}{subsection}

% Macros
%
\def\lim{\mathop{\mathrm{lim}}\nolimits}
\def\colim{\mathop{\mathrm{colim}}\nolimits}
\def\Spec{\mathop{\mathrm{Spec}}}
\def\Hom{\mathop{\mathrm{Hom}}\nolimits}
\def\Ext{\mathop{\mathrm{Ext}}\nolimits}
\def\SheafHom{\mathop{\mathcal{H}\!\mathit{om}}\nolimits}
\def\SheafExt{\mathop{\mathcal{E}\!\mathit{xt}}\nolimits}
\def\Sch{\mathit{Sch}}
\def\Mor{\mathop{\mathrm{Mor}}\nolimits}
\def\Ob{\mathop{\mathrm{Ob}}\nolimits}
\def\Sh{\mathop{\mathit{Sh}}\nolimits}
\def\NL{\mathop{N\!L}\nolimits}
\def\CH{\mathop{\mathrm{CH}}\nolimits}
\def\proetale{{pro\text{-}\acute{e}tale}}
\def\etale{{\acute{e}tale}}
\def\QCoh{\mathit{QCoh}}
\def\Ker{\mathop{\mathrm{Ker}}}
\def\Im{\mathop{\mathrm{Im}}}
\def\Coker{\mathop{\mathrm{Coker}}}
\def\Coim{\mathop{\mathrm{Coim}}}

% Boxtimes
%
\DeclareMathSymbol{\boxtimes}{\mathbin}{AMSa}{"02}

%
% Macros for moduli stacks/spaces
%
\def\QCohstack{\mathcal{QC}\!\mathit{oh}}
\def\Cohstack{\mathcal{C}\!\mathit{oh}}
\def\Spacesstack{\mathcal{S}\!\mathit{paces}}
\def\Quotfunctor{\mathrm{Quot}}
\def\Hilbfunctor{\mathrm{Hilb}}
\def\Curvesstack{\mathcal{C}\!\mathit{urves}}
\def\Polarizedstack{\mathcal{P}\!\mathit{olarized}}
\def\Complexesstack{\mathcal{C}\!\mathit{omplexes}}
% \Pic is the operator that assigns to X its picard group, usage \Pic(X)
% \Picardstack_{X/B} denotes the Picard stack of X over B
% \Picardfunctor_{X/B} denotes the Picard functor of X over B
\def\Pic{\mathop{\mathrm{Pic}}\nolimits}
\def\Picardstack{\mathcal{P}\!\mathit{ic}}
\def\Picardfunctor{\mathrm{Pic}}
\def\Deformationcategory{\mathcal{D}\!\mathit{ef}}
\usepackage{cleveref}
\crefname{originallemma}{original lemma}{original lemmas}
\Crefname{originallemma}{Original lemma}{Original lemmas}
\makeatletter
\newlabel{lemma-additivity-divisors-restricted}{{02QG}{source}{}{sourcepointer-1}{}}
\newlabel{lemma-additivity-divisors-restricted@cref}{{[originallemma][0][]02QG}{[1][0][]source}{}{}{}}
\makeatother
\crefname{originaldefinition}{original definition}{original definitions}
\Crefname{originaldefinition}{Original definition}{Original definitions}
\makeatletter
\newlabel{algebra-definition-ord}{{02MD}{source}{}{sourcepointer-2}{}}
\newlabel{algebra-definition-ord@cref}{{[originaldefinition][0][]02MD}{[1][0][]source}{}{}{}}
\makeatother
\title{Length difference for the ideal of denominators over a one dimensional Noetherian local ring without embedded primes}
\author{The Stacks Project Authors}
\date{stacks-project:chow.tex:local-snapshot-20261012; GNU FDL 1.2 or later}
\begin{document}
\maketitle
\noindent\textit{Editorial note.} All delivered mathematical text is the author's own, copied from the pinned Stacks Project chapter file chow.tex; no statement or proof step was rewritten, reordered or removed.\par\medskip
\noindent\textit{Editorial note.} Preamble provenance: the pinned chapter file opens with an input of the shared preamble and carries no preamble text of its own, so the project's single Stacks-wide preamble staged at sources/stacks-chow/preamble.tex is excerpted verbatim as two spans, lines 14-18 (Xy-pic setup) and lines 154-217 (the native theorem declarations and mathematical macros). Nothing was added, removed or redefined; the wrapper class and the amsmath/amssymb/amsthm/graphicx base packages are supplied by the assembler.\par\medskip
\noindent\textit{Editorial note.} The counted statement and proof contain no cross reference, no citation and no included file, so the delivered body is self contained.\par\medskip
\noindent\textit{Editorial note.} Canonical location of the counted proposition: Stacks tag 02QH (sources/stacks-chow/pinned-original-source.tex line 18594).\par\medskip
\noindent\textit{Allowed standard prerequisites (editorial).}\begin{itemize}
\item Additivity of divisors for the restricted situation with a and b. Same chapter, Lemma labelled lemma-additivity-divisors-restricted (source line 373), invoked at proof lines 18641-18643. Supplies the formula that turns length\_A(M/aM) - length\_A(M/bM) into the sum over the minimal primes; the reduction to that difference is carried out in the counted span.
\item The order function ord\_\{A/q\} on the residue field of a one dimensional quotient. Algebra, Definition labelled algebra-definition-ord (Stacks tag 02MD), invoked at proof line 18644. Fixes the meaning of the ord-functions appearing on the right hand side.
\end{itemize}\medskip
\ifdefined\hypertarget\hypertarget{sourcepointer-1}{}\fi
\noindent\textit{Source reference.} Original lemma 02QG: sources/stacks-chow/pinned-original-source.tex, source lines 373--373. Same-chapter pointer into the author environment lemma at source line 373; invoked at proof line 18642 of the counted proof. It is retained only as a navigation pointer; every load-bearing step is inside the delivered span.\par\medskip
\ifdefined\hypertarget\hypertarget{sourcepointer-2}{}\fi
\noindent\textit{Source reference.} Original definition 02MD: sources/stacks-chow/tags.txt, source lines 3440--3440. Cross-chapter standard prerequisite located in the project tag list staged in the local snapshot; invoked at proof line 18644 of the counted proof. It is retained only as a navigation pointer; every load-bearing step is inside the delivered span.\par\medskip
\begin{lemma}
\label{lemma-no-embedded-points-modules}
Let $A$ be a Noetherian local ring.
Let $M$ be a finite $A$-module.
Let $a, b \in A$.
Assume
\begin{enumerate}
\item $\dim(A) = 1$,
\item both $a$ and $b$ are nonzerodivisors in $A$,
\item $A$ has no embedded primes,
\item $M$ has no embedded associated primes,
\item $\text{Supp}(M) = \Spec(A)$.
\end{enumerate}
Let $I = \{x \in A \mid x(a/b) \in A\}$.
Let $\mathfrak q_1, \ldots, \mathfrak q_t$ be the minimal
primes of $A$. Then $(a/b)IM \subset M$ and
$$
\text{length}_A(M/(a/b)IM)
-
\text{length}_A(M/IM)
=
\sum\nolimits_i
\text{length}_{A_{\mathfrak q_i}}(M_{\mathfrak q_i})
\text{ord}_{A/\mathfrak q_i}(a/b)
$$
\end{lemma}

\begin{proof}
Since $M$ has no embedded associated primes, and since
the support of $M$ is $\Spec(A)$ we see that
$\text{Ass}(M) = \{\mathfrak q_1, \ldots, \mathfrak q_t\}$.
Hence $a$, $b$ are nonzerodivisors on $M$. Note that
\begin{align*}
& \text{length}_A(M/(a/b)IM) \\
& = \text{length}_A(bM/aIM) \\
& = \text{length}_A(M/aIM)
-
\text{length}_A(M/bM) \\
& = \text{length}_A(M/aM) + \text{length}_A(aM/aIM) - \text{length}_A(M/bM) \\
& = \text{length}_A(M/aM) + \text{length}_A(M/IM) - \text{length}_A(M/bM)
\end{align*}
as the injective map $b : M \to bM$ maps $(a/b)IM$ to $aIM$
and the injective map $a : M \to aM$ maps $IM$ to $aIM$.
Hence the left hand side of the equation of the lemma is
equal to
$$
\text{length}_A(M/aM) - \text{length}_A(M/bM).
$$
Applying the second formula of
Lemma \ref{lemma-additivity-divisors-restricted}
with $x = a, b$ respectively
and using Algebra, Definition \ref{algebra-definition-ord}
of the $\text{ord}$-functions we get the result.
\end{proof}

\end{document}
